\section{Arithmetic conductors and the exact modulus comparison}\label{sec:stark}
\noindent\textit{Source anchor A: preserved mathematical development; exact source path and hash are in the claim/source inventory. Retained local equation numbers carry this prefix.}\par
\hypertarget{A-fixed-conventions-and-inherited-arithmetic}{%
\subsection{Fixed conventions and inherited
arithmetic}\label{A-fixed-conventions-and-inherited-arithmetic}}

Set

\[
\omega=e^{2\pi i/3},\quad \beta=7^{1/3}\omega^2,\quad
J=2+3\omega,\quad \bar J=-1-3\omega,\quad J\bar J=7.
\]

Write

\[
F_0=\mathbb Q(\omega),\qquad K=\mathbb Q(\beta),\qquad
\epsilon=2-\beta,\qquad \mathcal R=-\log(2-7^{1/3})>0.
\]

Here \(\mathcal R\) is the regulator of \(K\), and

\[
\log|\epsilon|=\mathcal R/2. \tag{A1.1}
\]

Indeed, \(N_{K/\mathbb Q}(\epsilon)=1\) and its real image is
\(2-7^{1/3}\). The class number is three and \(\epsilon\) is a
fundamental unit. These arithmetic facts are recorded in BCG, section
5.2; the earlier unramified-cubic note also proves the class-number and
ray-field identification used here.

Take

\[
\mathfrak f=(3,\beta-1),\qquad
E=K(\omega,\theta),\qquad \theta=2\cos(6\pi/7).
\]

The inherited arithmetic proof identifies \(E\) as the narrow ray class
field of modulus \(\mathfrak f\) and the real place. It has degree six
over \(K\) and degree eighteen over \(\mathbb Q\).

Let

\[
X=\frac1\beta\sum_{a=1}^{6}\chi_7(a)e^{6\pi ia/7},
\]

where the character values for \(a=1,\ldots,6\) are
\(1,\omega,\omega^2,\omega^2,\omega,1\). Then

\[
X^3=J,\qquad E=\mathbb Q(\beta,X,\omega). \tag{A1.2}
\]

This fixes a specific cube-root branch. The preceding CM note proves,
with this branch,

\[
X=\frac{\Omega^2}{\wp_{\Lambda_0}(1/J)},\qquad
\Lambda_0=\mathbb Z+\mathbb Z\omega,\qquad
\Omega=e^{\pi i/6}\frac{\Gamma_{\rm Euler}(1/3)^3}{2\pi}. \tag{A1.3}
\]

That proof uses actual seven-torsion and a lattice estimate for the
branch, not numerical recognition.

The arithmetic Frobenius \(\sigma\) at the norm-two prime acts by

\[
\sigma(\beta)=\beta,\quad \sigma(\omega)=\omega^2,\quad
\sigma(X)=\omega\beta/X. \tag{A1.4}
\]

The norm-five prime has the same Artin class. In particular, \(\sigma\)
generates \(\operatorname{Gal}(E/K)\).

\textbf{Actual complex conjugation is a different automorphism.} Denote
it by \(\kappa\). In the chosen embedding,

\[
\kappa(\beta)=\omega^2\beta,\quad \kappa(\omega)=\omega^2,
\quad \kappa(X)=\omega\beta/X. \tag{A1.5}
\]

For the last formula, the conjugate Gauss resolvent is \(7/G\), and
\(\bar\beta=\omega^2\beta\). Formula (A1.5) follows directly. It must
not be replaced by the involution \(\sigma^3\), which fixes \(K\).

\hypertarget{A-the-major-change-of-route-artin-induction}{%
\subsection{The major change of route: Artin
induction}\label{A-the-major-change-of-route-artin-induction}}

The splitting field \(F_0(\beta)/\mathbb Q\) has group \(S_3\), while
\(\mathbb Q(\theta)/\mathbb Q\) is cyclic cubic. Their intersection is
Galois over \(\mathbb Q\); it cannot have degree three because \(S_3\)
has no cyclic cubic quotient. Hence

\[
\operatorname{Gal}(E/\mathbb Q)=S_3\times C_3. \tag{A2.1}
\]

Use generators

\[
r:\beta\mapsto\omega\beta,\qquad
c:\omega\mapsto\omega^2,\qquad
s_0:\theta\mapsto\theta^2-2,
\]

where each generator fixes the other specified primitive generators as
appropriate. Thus \(r,c\) generate the \(S_3\) factor, \(c\) fixes
\(\beta,\theta\), and \(s_0\) generates the second factor. We have

\[
\operatorname{Gal}(E/K)=\langle c,s_0\rangle,
\qquad \sigma=cs_0,
\qquad \operatorname{Gal}(E/F_0)=\langle r,s_0\rangle.
\]

Let \(\varrho(r)=\omega,\varrho(s_0)=1\), and let
\(\eta(s_0)=\omega,\eta(r)=1\). The three odd ray characters of \(K\)
are

\[
\chi_j(c)=-1,\qquad \chi_j(s_0)=\omega^j,\qquad j=0,1,2.
\]

The elementary character identity

\[
\operatorname{Ind}_{C_2}^{S_3}(\operatorname{sgn})
=\operatorname{sgn}\oplus\operatorname{Std},
\qquad
\operatorname{Std}=\operatorname{Ind}_{C_3}^{S_3}(\varrho)
\]

and induction invariance of Artin L-functions give

\[
\boxed{L_K(s,\chi_j)
=L_{\mathbb Q}(s,\chi_{-3}\chi_7^j)\,
L_{F_0}(s,\varrho\eta^j).} \tag{A2.2}
\]

For \(j=0\), use the primitive Dirichlet character of conductor three,
not its imprimitive extension to modulus twenty-one. All three odd ray
characters are ramified at \(\mathfrak f\), so their Artin L-functions
agree with the partial-zeta convention omitting that prime. Induction
includes the local factors at ramified primes through the
inertia-invariant spaces. No unproved Artin holomorphy statement is
used.

The even characters make no contribution to the derivative at zero. The
nontrivial even characters have at least two archimedean zeros. The
trivial character has the zero of \(\zeta_K\), together with the extra
zero of \(1-3^{-s}\). This also agrees with BCG, section 3.1.

Exact finite Dirichlet sums give

\[
D_0=L(0,\chi_{-3})=1/3,\qquad
D_1=L(0,\chi_{-3}\chi_7)=-2\omega,\qquad
D_2=-2\omega^2. \tag{A2.3}
\]

For example, the latter two use
\(-21^{-1}\sum_{a=1}^{21}a\chi_{-3}(a)\chi_7(a)^j\).

For the untwisted cubic factor,

\[
L_{F_0}(s,\varrho)=\zeta_K(s)/\zeta_{\mathbb Q}(s).
\]

The analytic class number formula, \(h_K=3\), and
\(\zeta_{\mathbb Q}(0)=-1/2\) give

\[
L_{F_0}'(0,\varrho)=3\mathcal R. \tag{A2.4}
\]

Thus the cubic regulator contributes a known normalization; it is not a
free scalar left over from a cocycle.

\hypertarget{A-a-six-class-imaginary-quadratic-calculation}{%
\subsection{A six-class imaginary-quadratic
calculation}\label{A-a-six-class-imaginary-quadratic-calculation}}

Put \(\lambda=\varrho\eta\). Since both \(r\) and \(s_0\) send \(X\) to
\(\omega^2X\), \(\lambda\) is the inverse Kummer character for
\(F_0(X)/F_0\). Its conductor is

\[
\mathfrak m=3J\mathscr O_{F_0}. \tag{A3.1}
\]

Here is an explicit verification. The Eisenstein integers are principal,
and the six units map bijectively to \((\mathscr O_{F_0}/3)^\times\).
Every ideal coprime to three therefore has a unique generator
\(\alpha\equiv1\pmod3\). Since \(J\equiv-1\pmod3\), cubic reciprocity
gives

\[
\lambda((\alpha))=(J/(\alpha))_3^2=(\alpha/J)_3^2. \tag{A3.2}
\]

This proves that the conductor divides \(3J\). It must contain \(J\),
because the valuation of the Kummer radicand there is one. If its
exponent at the prime above three were at most one, its conductor would
divide \((1-\omega)J\). The ray class group of that modulus has order
two, so it has no nontrivial cubic character. Hence the conductor is
exactly (A3.1).

The ray group modulo \(3J\) has order

\[
\frac{63(1-1/3)(1-1/7)}6=6.
\]

Its six classes have unique primary rational representatives

\[
1,4,10,13,16,19.
\]

Reduction modulo \(J\) sends \(\omega\) to \(4\in\mathbb F_7\). Formula
(A3.2) groups the representatives as follows:

\begin{longtable}[]{@{}lll@{}}
\toprule\noalign{}
Index \(m\) & Representatives & \(\lambda\)-value \\
\midrule\noalign{}
\endhead
\bottomrule\noalign{}
\endlastfoot
0 & 1, 13 & 1 \\
1 & 4, 10 & \(\omega\) \\
2 & 16, 19 & \(\omega^2\) \\
\end{longtable}

Let \(Z_m(s)\) be the sum of the two ray partial zeta functions in row
\(m\), and write \(\ell_m=Z_m'(0)\). Then

\[
L_{F_0}'(0,\lambda)=\ell_0+\omega\ell_1+\omega^2\ell_2,
\qquad
L_{F_0}'(0,\lambda^2)=\ell_0+\omega^2\ell_1+\omega\ell_2. \tag{A3.3}
\]

The sum of all six ray partial zeta functions is

\[
\zeta_{F_0}(s)(1-3^{-s})(1-7^{-s}),
\]

whose derivative at zero vanishes. Therefore

\[
\ell_0+\ell_1+\ell_2=0. \tag{A3.4}
\]

Conjugation of the base field by \(c\) sends \(\varrho\eta^2\) to
\(\varrho^{-1}\eta^2=\lambda^2\), so

\[
L_{F_0}(s,\varrho\eta^2)=L_{F_0}(s,\lambda^2).
\]

Combining (A2.2)--(A3.4) and performing inverse Fourier transformation
on \(C_2\times C_3\) yields

\[
\boxed{\zeta_{\mathfrak f}'(c^e s_0^k,0)
=(-1)^e\left(\frac{\mathcal R}{6}-\ell_{k-1}\right).} \tag{A3.5}
\]

All subscripts on \(\ell\) are taken modulo three. To check the
constants directly, the two nontrivial Fourier contributions are

\[
-\frac13\left(\omega^{1-k}L_{F_0}'(0,\lambda)
+\omega^{2-2k}L_{F_0}'(0,\lambda^2)\right).
\]

In (A3.3), the coefficient of \(\ell_{k-1}\) in the parentheses is two
and the other two coefficients are minus one. Equation (A3.4) turns the
parentheses into \(3\ell_{k-1}\).

\textbf{Euler-factor warning.} Replacing \(3J\) by \(21=3J\bar J\) would
remove the Euler factor at the other prime over seven. Its
\(\lambda\)-value is one, so this inserts \(1-7^{-s}\) and destroys the
desired first derivative. That imprimitive derivative is zero. The
checker verifies this as a separate normalization control, not as the
desired L-value.

\hypertarget{A-kronecker-normalization-fixed-explicitly}{%
\subsection{Kronecker normalization fixed
explicitly}\label{A-kronecker-normalization-fixed-explicitly}}

For \(z=a\omega+b\notin\Lambda_0\), define the positive quantity

\[
H(z)=e^{-\pi\operatorname{Im}(\omega)B_2(a)}
|1-e^{2\pi iz}|
\prod_{n\ge1}|(1-e^{2\pi i(n\omega+z)})(1-e^{2\pi i(n\omega-z)})|,
\tag{A4.1}
\]

where initially \(0\le a<1\) and \(B_2(a)=a^2-a+1/6\). This is the
absolute value of the usual Siegel function. Its value is independent of
the representative of \(z\) modulo \(\Lambda_0\).

For the shifted Epstein series

\[
\mathcal E_z(s)=\sum_{m,n\in\mathbb Z}|m+n\omega+z|^{-2s},
\]

continued from \(\Re s>1\), the precise normalization is

\[
\mathcal E_z(0)=0,\qquad \mathcal E_z'(0)=-2\log H(z). \tag{A4.2}
\]

A normalization derivation is given in the normalization subsection below (source A8). In particular, no
factor of two is inferred from an unspecified Mellin convention.

For the ray class represented by \(\alpha\), all its integral ideals are
generated exactly once by the elements of \(\alpha+3J\mathscr O_{F_0}\).
The congruence condition eliminates all nontrivial units. Therefore its
partial zeta function is

\[
63^{-s}\mathcal E_{\alpha/(3J)}(s).
\]

There is no scale contribution to its derivative, because
\(\mathcal E_z(0)=0\). Set

\[
V_0=H(1/(3J))H(13/(3J)),
\] \[
V_1=H(4/(3J))H(10/(3J)),\qquad
V_2=H(16/(3J))H(19/(3J)).
\]

Then

\[
\ell_m=-2\log V_m,
\qquad V_0V_1V_2=1. \tag{A4.3}
\]

\hypertarget{A-the-cm-products-collapse-to-three-explicit-units}{%
\subsection{The CM products collapse to three explicit
units}\label{A-the-cm-products-collapse-to-three-explicit-units}}

Let \(u=1/J\) and \(T=\bar J/3\equiv2/3\pmod{\Lambda_0}\). The three
pairs of points above are, respectively,

\[
\{T\pm2u\},\quad\{T\pm u\},\quad\{T\pm3u\}. \tag{A5.1}
\]

Indeed, \(T=7u/3\), so their numerators over \(3J\) are \(1,13\),
\(4,10\), and \(-2,16\). The last numerator \(-2\) may be replaced by
\(19\) since \(21/(3J)=\bar J\in\Lambda_0\).

On the normalized curve \(E_0:y^2+y=x^3\), set

\[
x(z)=\wp_{\Lambda_0}(z)/\Omega^2.
\]

The previous seven-torsion proof gives

\[
x(u)=X^{-1},\quad x(2u)=\omega^2X^{-1},\quad x(3u)=\omega X^{-1}. \tag{A5.2}
\]

Also,

\[
x(T)=-\omega. \tag{A5.3}
\]

For completeness, \(T\) is nonzero three-torsion and its normalized
x-coordinate is a root of \(3x(x^3+1)\). The lattice is invariant under
ordinary conjugation, so \(\wp(1/3)\) and its derivative are real. Since
\(g_3(\Lambda_0)>0\), the real differential equation
\(\wp'^2=4\wp^3-g_3\) implies \(\wp(1/3)>0\). The argument of
\(\Omega^2\) is \(\pi/3\), selecting \(-\omega\) among the three roots
of \(x^3=-1\).

The classical sigma addition identity is

\[
\frac{\sigma_\Lambda(T+z)\sigma_\Lambda(T-z)}{
\sigma_\Lambda(T)^2\sigma_\Lambda(z)^2}
=\wp_\Lambda(z)-\wp_\Lambda(T). \tag{A5.4}
\]

The theta expression for (A4.1), or the sigma--theta relation, turns
(A5.4) into the same identity for absolute Green functions up to a
constant independent of \(z\). The quadratic exponential corrections
cancel by the parallelogram identity. Furthermore,
\(H(\omega z)=H(z)=H(-z)\), either from (A4.2) and the isometries of
\(\Lambda_0\), or from the theta transformations.

Since \(2u\equiv\omega^2u\) and \(3u\equiv-\omega u\), all three factors
\(H(ku)\) are equal. Thus for one positive constant \(C\), the same for
all three rows,

\[
V_0=C|\omega^2/X+\omega|,\quad
V_1=C|1/X+\omega|,\quad
V_2=C|\omega/X+\omega|. \tag{A5.5}
\]

Multiplication and (A4.3) give \(C^3|X^{-3}+1|=1\). Cubing each equation
in (A5.5) now removes the unknown constant exactly:

\[
\boxed{V_m^3=|a_m|,} \tag{A5.6}
\]

where the three algebraic values simplify to

\[
\begin{aligned}
a_0&=\frac{(\omega^2+\omega X)^3}{1+J}
=(1+\omega)X^2-X+1,\\
a_1&=\frac{(1+\omega X)^3}{1+J}
=-X^2+(1+\omega)X+1,\\
a_2&=\frac{(\omega+\omega X)^3}{1+J}
=-\omega X^2-\omega X+1.
\end{aligned} \tag{A5.7}
\]

Exact reduction gives

\[
a_0+a_1+a_2=3,\qquad
\sum_{i<j}a_i a_j=6\omega,\qquad a_0a_1a_2=1.
\]

In particular, the \(a_m\) are algebraic units, satisfying

\[
a_m^3-3a_m^2+6\omega a_m-1=0. \tag{A5.8}
\]

Equations (A4.3) and (A5.6) prove

\[
\boxed{\ell_m=-\frac23\log|a_m|.} \tag{A5.9}
\]

This is the evaluated comparison. The unknown logarithmic constants have
been replaced by explicit algebraic units without using any
elliptic-gamma algebraicity assumption.

Substitution into (A3.5) yields the explicit partial-zeta theorem

\[
\boxed{\zeta_{\mathfrak f}'(c^e s_0^k,0)
=(-1)^e\left(\frac{\mathcal R}{6}+\frac23\log|a_{k-1}|\right).} \tag{A5.10}
\]

\hypertarget{A-identifying-the-previously-constructed-w}{%
\subsection{Identifying the previously constructed
W}\label{A-identifying-the-previously-constructed-w}}

Recall the integral cubic coordinate

\[
\nu=(X+\beta/X-\beta^2)/3,\quad v=-\beta^2\nu-2\nu^2,
\qquad v^3-\beta v^2+1=0.
\]

The earlier unit is

\[
W=(A_W+\omega B_W)/3,
\]

\[
\begin{aligned}
A_W&=(21\beta-42)v^2+(-15\beta^2+30\beta)v-\beta^2-4\beta+11,\\
B_W&=3\beta^2v^2+(18\beta-42)v-11\beta^2+13\beta+7.
\end{aligned} \tag{A6.1}
\]

Let \(N_\kappa(z)=z\kappa(z)=|z|^2\) in the fixed embedding. The six
exact identities needed for the comparison are

\[
\boxed{N_\kappa(\sigma^hW)
=N_\kappa(\epsilon a_{h+1}^2)^{(-1)^h},\quad 0\le h<6.} \tag{A6.2}
\]

These are identities in the explicitly presented degree-eighteen field,
not fits to the logarithms. They can be checked by substituting (A1.4),
(A1.5), (A5.7), and (A6.1), then reducing using

\[
\beta^3=7,\quad X^3=2+3\omega,\quad \omega^2+\omega+1=0.
\]

A basis for the rational calculation is \(\beta^iX^j\omega^k\),
\(0\le i,j<3\), \(0\le k<2\). For odd \(h\), multiply the two norms
instead of dividing: the residual is their product minus one. Both
supplied implementations obtain the exact zero polynomial in all six
cases.

Taking real logarithms of positive norms gives

\[
\log|\sigma^h W|=(-1)^h\left(\frac{\mathcal R}{2}+2\log|a_{h+1}|\right)
=-3\zeta_{\mathfrak f}'(\sigma^{h+5},0). \tag{A6.3}
\]

Consequently \(W_{\rm St}=\sigma(W)\) satisfies

\[
\zeta_{\mathfrak f}'(\sigma^h,0)
=-\frac13\log|\sigma^hW_{\rm St}|
=-\frac16\log|\sigma^hW_{\rm St}|^2. \tag{A6.4}
\]

The one-step shift is essential. The unshifted \(W\) belongs to the
original raw gamma value's ray-class indexing, not the identity class.

\hypertarget{A-completing-the-stark-conditions-for-this-extension}{%
\subsubsection{Completing the Stark conditions for this
extension}\label{A-completing-the-stark-conditions-for-this-extension}}

The earlier note proves that \(W\) is an algebraic unit,
\(\mu(E)=\mu_6\), all its embeddings above the real place of \(K\) have
modulus one, and adjoining a sixth root produces an abelian extension of
\(K\). These statements are preserved by \(\sigma\). Together with
(A6.4), this verifies the three conditions stated by BCG in section 3.3
for this specific field and modulus.

The archived source checks rechecked the following inherited algebraic certificates. In
particular,

\[
W^5=\sigma(W)U_{\rm alg}^6,
\]

and, for \(c_6=-\omega^2\),

\[
W^{21}=c_6U_{\rm alg}^{25}\sigma(U_{\rm alg})^5\sigma^2(U_{\rm alg}).
\]

For \(t^6=W\), the lift \(\varrho(t)=t^5/U_{\rm alg}\) commutes with
\(t\mapsto c_6t\), and satisfies \(\varrho^6(t)=c_6^4t\). Irreducibility
is checked at the residue assignment \((\beta,v,\omega)=(32,99,48)\)
over \(\mathbb F_{181}\), where \(W=23\) and \(T^6-23\) is irreducible.
The earlier local-integrality argument validates this reduction. The
resulting abelian group is \(C_2\times C_{18}\).

These arithmetic statements concern \(W\), not the still unproved
algebraicity of the gamma value.

\hypertarget{A-closing-the-actual-gamma-modulus-gap}{%
\subsection{Closing the actual gamma modulus
gap}\label{A-closing-the-actual-gamma-modulus-gap}}

The original analytic value is

\[
U_\Gamma=\frac{\Gamma_{\rm ell}(1/3;\tau,s_{\rm per})^5}
{\Gamma_{\rm ell}(5/3;5\tau,5s_{\rm per})},
\quad
\tau=(\beta^2+2\beta)/15,\quad s_{\rm per}=-(\beta+2)/15.
\]

BCG section 5.2 identifies this raw expression with the geometric value
at ray-class index five. The norm-five smoothing ideal has Artin class
\(\sigma\). Their Theorem 3.2 therefore gives, unconditionally,

\[
\log|U_\Gamma|^2
=\zeta_{\mathfrak f}'(1,0)-5\zeta_{\mathfrak f}'(\sigma^5,0). \tag{A7.1}
\]

Using (A6.3),

\[
\zeta_{\mathfrak f}'(1,0)-5\zeta_{\mathfrak f}'(\sigma^5,0)
=\frac13\left(5\log|W|-\log|\sigma(W)|\right).
\]

The exact algebraic smoothing identity \(W^5/\sigma(W)=U_{\rm alg}^6\)
turns this into \(2\log|U_{\rm alg}|\). Thus

\[
\boxed{|U_\Gamma|=|U_{\rm alg}|.} \tag{A7.2}
\]

There is also a smaller explicit algebraic expression:

\[
\log|U_\Gamma|^2
=\mathcal R-\ell_2-5\ell_1
=\mathcal R+\frac23\log|a_2a_1^5|,
\]

or

\[
\boxed{|U_\Gamma|^3=|\epsilon^3a_2a_1^5|.} \tag{A7.3}
\]

In particular the positive real number \(|U_\Gamma|\) is algebraic. This
conclusion does not extend automatically to the complex number
\(U_\Gamma\).

\hypertarget{A-the-whole-logarithmic-ray-packet}{%
\subsubsection{The whole logarithmic ray
packet}\label{A-the-whole-logarithmic-ray-packet}}

The same argument evaluates the moduli of the admissible geometric gamma
values for every ray class and every smoothing ideal covered by BCG's
theorem:

\[
|\mathbf\Gamma_{\mathfrak a,h}(\mathfrak f\mathfrak b^{-1})|^6
=\left|(\sigma_{\mathfrak b}W_{\rm St})^{N(\mathfrak a)-\sigma_{\mathfrak a}}\right|.
\tag{A7.4}
\]

This is an equality of absolute values. It is not an Artin action on an
already algebraic gamma value and does not specify the phase of a sixth
root.

\hypertarget{A-derivation-of-the-shifted-epstein-normalization}{%
\subsection{Derivation of the shifted Epstein
normalization}\label{A-derivation-of-the-shifted-epstein-normalization}}

Let \(\tau=x+iy\) with \(y>0\), and \(z=a\tau+b\), initially \(0<a<1\).
Apply Poisson summation in the integer coordinate \(m\) of

\[
\mathcal E_z(s)=\sum_{n,m\in\mathbb Z}
\left((m+b+(n+a)x)^2+(n+a)^2y^2\right)^{-s}.
\]

The zero Fourier mode is

\[
\sqrt\pi\,\frac{\Gamma(s-1/2)}{\Gamma(s)}\,
 y^{1-2s}\bigl[\zeta(2s-1,a)+\zeta(2s-1,1-a)\bigr].
\]

At zero its value vanishes. Its derivative equals \(2\pi y B_2(a)\),
using \(\Gamma(-1/2)=-2\sqrt\pi\), \(1/\Gamma(s)=s+O(s^2)\), and
\(\zeta(-1,a)=-B_2(a)/2\).

The nonzero Fourier modes, obtained from the usual Gaussian integral,
are exponentially convergent near \(s=0\) for \(0<a<1\). Their
derivative sums to

\[
-2\log|1-e^{2\pi iz}|
-2\sum_{n\ge1}\log|(1-e^{2\pi i(n\tau+z)})(1-e^{2\pi i(n\tau-z)})|.
\]

Combining the two terms gives (A4.2). Extension to \(a=0\), away from
lattice points, follows by continuation in \(z\). Equivalently,

\[
H_\tau(z)=e^{-\pi(\Im z)^2/y}
\left|\frac{\theta_1(\pi z;e^{\pi i\tau})}{\eta_{\rm Ded}(\tau)}\right|.
\]

This representation and the sigma--theta relation justify the absolute
CM-product addition step above (source A5), including the cancellation of quadratic
exponential corrections.
